\documentclass[11pt]{article}
\usepackage[T1]{fontenc}
\usepackage[utf8]{inputenc}
\usepackage{lmodern,amsmath,amssymb,amsthm,booktabs,array,microtype}
\usepackage[margin=26mm]{geometry}
\usepackage[hidelinks]{hyperref}
\usepackage{xurl}
\newtheorem{theorem}{Theorem}
\newtheorem{lemma}[theorem]{Lemma}
\newtheorem{proposition}[theorem]{Proposition}
\theoremstyle{remark}\newtheorem*{remark}{Scope}
\newcommand{\T}{\mathbb T}
\newcommand{\eps}{\varepsilon}
\newcommand{\poly}{\operatorname{poly}}
\newcommand{\lgof}[1]{\log_2\!\left(#1\right)}
\allowdisplaybreaks
\setlength{\emergencystretch}{2em}
\title{A conditional Subset Sum refinement to exponent 0.4876723}
\author{Yanni Bills}
\date{21 October 2026}
\begin{document}
\maketitle
\begin{abstract}
We derive a conditional $O_c(2^{0.4876723n})$ bound by modifying the transfer of existing checksum certificates in the claimed
Subset Sum construction of Kapila, following Wharton's conditional parameter
refinement. A certificate whose lower-region multiplier is zero remains
usable when that region's lower boundary moves. For filter~3 this allows
its low-region certificate to stop at the new sampling boundary, replacing
an expensive uniform-transfer loss by a smaller two-constant shift.
We prove the domination and full-circle integral identities, preserving all
filters and requiring no new quadrature. The resulting algorithmic claim
remains conditional on explicitly identified inherited algorithmic claims
and continuous moment inequalities. The internal total work exponent is $0.48767228002$, strictly below the advertised exponent after all stated accuracy costs. Exact rational interval checks establish the new scalar inequalities. No end-to-end Lean or site-verification claim is made.
\end{abstract}

\section{Problem, antecedents and the conditional claim}
The problem is to decide, for positive integers $a_1,\ldots,a_n$ and
nonnegative $t$, whether $\sum_{i\in I}a_i=t$ for some $I\subseteq[n]$.
Repeated values represent distinct indices; $I=\varnothing$ is permitted.
We use the live classical randomized word-RAM specification~\cite{benchmark}:
for $n\ge2$, $b$ is the largest input or target bit length and words have
$\lceil4(n+b+\log_2(n+2))\rceil$ bits. Overflow is charged through multiword
operations, all preparation counts, retained randomness occupies storage,
and there is no advice or external table. One program, independent of $c$,
must halt on every run, succeed with probability at least $2/3$ on each
input, and have the stated operation bound on every random outcome for
every fixed polynomial-bit regime $b\le n^c$.

Write $\mathcal K$ for Kapila's pinned manuscript~\cite{kapila} and
$\mathcal W$ for Wharton's leaderboard-linked note~\cite{wharton}.
The latter refines internal work exponent $0.48769$ to $0.48767595$
and advertises the conditional exponent $0.487676$; its original
moment bounds remain hypotheses. Both inherit the OpenAI framework~\cite{openai}.
Our change is the special boundary rule proved in Section~\ref{sec:new},
not merely choosing a smaller advertised exponent from existing slack.
The source comparison is bounded: we checked the pinned papers and the
current Wharton and Kapila default-branch trees on 21 October 2026,
at revisions \texttt{bed8e18} and \texttt{44c3496}, respectively,
including the related transfer, finite-data and verification files.
Those applications retain the uniform low-region transfer or the original
region rules; we did not identify this particular filter-3 boundary reuse.
This establishes the distinction from the inspected sources, rather than
a comprehensive claim of priority. The exact file and hash record is
\texttt{source-audit/novelty-check.txt} and its accompanying inventory.
The unchanged algorithms, filters, moment data and earlier transfer ideas
are credited to those works.

\paragraph{Inherited hypotheses.}
The conditional result uses the isolation, deficiency and finite
implementation claims in Sections~2 and~6 of~\cite{openai}; the
averaged-layer, preliminary-search and alias arguments in
Sections~3--5 of~$\mathcal K$; and its block-table and sampled-checksum
analysis in Section~6, with the parameter dependence documented in
Section~\ref{sec:admissible}. It also assumes the following precise
continuous moment inequalities for the original 430 certificates:
\begin{equation}\label{eq:oldmom}
 \lgof{I_{v,i}}<\eta_b-m_i,\qquad
 (m_0,m_1,m_{2:8},m_9)=(0.00025,0.00025,4,17),
\end{equation}
where $\eta_b=0.51613$ and the integrals are defined below.
These are inherited claimed inequalities, not integrals established by
our scalar checker. We do not assume the new boundary-transfer lemma
or our new scalar feasibility conclusions: they require the proofs and
exact certificate in this note. Any unresolved new numerical step prevents
a final exponent claim.

\begin{theorem}[Conditional refinement]\label{thm:main}
Assume the inherited algorithmic claims specified above and the original
moment inequalities~\eqref{eq:oldmom}. The fixed rational construction
parameters in Section~\ref{sec:admissible}, with the filter-3 certificate
boundary rule of Lemma~\ref{lem:new}, yield a single randomized classical
word-RAM program that halts on every input and is correct with probability
at least $2/3$. For each fixed positive integer $c$, there are $C_c,N_c>0$
such that every random execution on $n\ge N_c$, $b\le n^c$, uses at most
$C_c2^{0.4876723n}$ word operations in the stated model. The program does
not depend on $c$.
\end{theorem}
The proof combines Lemmas~\ref{lem:old}--\ref{lem:new}, the exact
feasibility proof in Proposition~\ref{prop:numeric}, and the
algorithmic/absorption argument in Proposition~\ref{prop:algorithm}.
This is a new conditional analysis within the existing construction,
not a proof that the inherited premises have been verified.

\section{Scores, certificates and the existing transfer}
All logarithms in scores are base two; $\T=\mathbb R/(2\pi\mathbb Z)$
has normalized measure. The fixed 43 filters define scores
$u_v=A_v+E_v-D_v$ and $w_v=2D_v-A_v-E_v$ from positive densities,
as in Section~3 of~\cite{openai} and Section~6.1 of~$\mathcal K$.
For a pattern, denote its approximate scores by $U_v,W_v$ and select the
least $v$ maximizing $U_v$. Thus $U_v-U_j\ge0$ for every comparison row.
The relative-entropy order and its prescribed rounding give $W_v\ge0$.
Write $U,W$ for the selected scores and put
\begin{equation}\label{eq:budgets}
 \Delta=\eta-\tau>0,\qquad\Gamma=2\tau-\eta>0,\qquad
 \alpha=\begin{cases}1,&W\le\Gamma,\\ \Gamma/W,&W>\Gamma.\end{cases}
\end{equation}
Only patterns satisfying $\alpha U<\Delta$ are sampled.
The old budgets are $\Delta_b=0.02844$ and $\Gamma_b=0.45925$.
Throughout the transfer assume
\begin{equation}\label{eq:qK}
 q=\Gamma/\Gamma_b\in(0,1),\quad
 A=\Delta/\Gamma,\quad A_b=\Delta_b/\Gamma_b,\quad
 K=\Delta/q-\Delta_b>0.
\end{equation}
In particular $A-A_b=K/\Gamma_b>0$.

For filter $v$ and old region $i$, let $z\ge0$ be its multiplier vector
and $G_b(x)$ its affine row vector. The old integral is
\begin{equation}\label{eq:integral}
 I_{v,i}=\int_{\T}(2+2\cos x)\,2^{z\cdot G_b(x)}\,dx.
\end{equation}
For the low region $W\le\Gamma_b$, the first rows are
$(\Delta_b-u_v,\Gamma_b-w_v)$ and $z_0\ge1$.
A bounded old band $\Gamma_b L<W\le\Gamma_b H$ has rows
$(A_bw_v-u_v,w_v-\Gamma_bL,\Gamma_bH-w_v)$ and $z_0L\ge1$.
The tail $W>15\Gamma_b$ has the first two of those rows, with $L=15$.
The remaining rows are selected comparisons $u_v-u_j$.
The old boundaries are
$1,1.25,1.6,2,2.6,3.5,5,8,15,\infty$, multiplied by $\Gamma_b$.

\begin{lemma}[Existing transfer, with domination]\label{lem:old}
For unchanged old regions the following modifications give ideal pattern-level
domination $-\alpha U\le-\Delta+z'\cdot G'^\pi$:
\begin{enumerate}
\item In the low region replace $\Delta_b-u_v$ by $\Delta/q-u_v$.
The integral becomes $2^{z_0K}I_{v,0}$.
\item In a bounded band replace $A_bw_v-u_v$ by $Aw_v-u_v$ and
increase the upper-row multiplier by $z_0(A-A_b)$. The integral becomes
$2^{z_0KH}I_{v,i}$.
\item In the tail replace $A_b$ by $A$ with unchanged multipliers.
If $w_v\le B$ pointwise, the integral is at most
$2^{Bz_0K/\Gamma_b}I_{v,9}$.
\end{enumerate}
Here $G'^\pi$ denotes rows evaluated at pattern scores.
\end{lemma}
\begin{proof}
This is the transfer of~$\mathcal W$, included to specify every region
used in our modification. For $W\le\Gamma_b$, $q\le\alpha\le1$.
Sampling gives $U<\Delta/\alpha\le\Delta/q$. Hence
\[
 -\alpha U=-\Delta+\alpha(\Delta/\alpha-U)
 \le-\Delta+z_0(\Delta/q-U).
\]
All other pattern rows are nonnegative. The pointwise exponent changes
by $z_0(\Delta/q-\Delta_b)=z_0K$.
In a bounded band, $\alpha<q/L<1/L\le z_0$ and $\alpha AW=\Delta$,
so sampling and $AW-U>0$ give
\[
 -\alpha U=-\Delta+\alpha(AW-U)\le-\Delta+z_0(AW-U).
\]
All remaining rows have nonnegative pattern values. The new upper
multiplier cancels the changed $w_v$ term:
\[
 z_0(A-A_b)w_v+z_0(A-A_b)(\Gamma_bH-w_v)=z_0KH.
\]
The tail has $\alpha<q/15\le z_0$ and the same domination. Its exponent
changes by $z_0(A-A_b)w_v\le Bz_0K/\Gamma_b$. Integrands are nonnegative,
so integration preserves the asserted comparisons.
\end{proof}

\section{The moved boundary for filter 3}\label{sec:new}
Appendix~C of~$\mathcal K$ uses zero-based filter indices. For filter~3
its low certificate has comparisons $\{5,16\}$ and multipliers
\[
 z^{3,0}=(17.211,27.1977,8.3302,3.0053).
\]
Its first bounded-band certificate has comparisons $\{5,8,16,35\}$ and
\[
 z^{3,1}=(7.7547,0,0,4.2699,2.2635,0.461,3.4032).
\]
These terminating decimals are exact rational source data. In particular,
the multiplier of $w_3-\Gamma_b$ is exactly zero, not approximately zero.

\begin{lemma}[Movable lower boundary]\label{lem:new}
Use filter~3's modified low certificate only for $W\le\Gamma$,
replacing both constant rows by $(\Delta-u_3,\Gamma-w_3)$.
Use its old first-band certificate throughout
$\Gamma<W\le1.25\Gamma_b$, replacing $A_b$ by $A$ and increasing
its upper-row multiplier by $7.7547(A-A_b)$.
Both certificates give ideal domination for every sampled pattern in
their new regions. Their respective integrals are exactly
\begin{align}
 I'_{3,0}&=2^{\Lambda_0}I_{3,0},&
 \Lambda_0&=17.211(\Delta-\Delta_b)+27.1977(\Gamma-\Gamma_b),\label{eq:speciallow}\\
 I'_{3,1}&=2^{\Lambda_1}I_{3,1},&
 \Lambda_1&=7.7547(1.25)K.\label{eq:specialband}
\end{align}
These are full-circle integral identities, not integrals restricted to
patterns satisfying the region inequalities.
\end{lemma}
\begin{proof}
In the new low region $\alpha=1$ and sampling implies $U<\Delta$.
Since $17.211\ge1$, and since the new upper-region row and all
comparison rows have nonnegative pattern values,
\[
 -U=-\Delta+(\Delta-U)
 \le-\Delta+17.211(\Delta-U)\le-\Delta+z^{3,0}\cdot G'^\pi.
\]
Changing the two constants changes the exponent pointwise by exactly
$\Lambda_0$, proving~\eqref{eq:speciallow}. Its sign need not be positive.

In the extended band $W>\Gamma$, so $\alpha=\Gamma/W<1<7.7547$.
Sampling gives $AW-U>0$ and $\alpha AW=\Delta$, whence
\[
 -\alpha U\le-\Delta+7.7547(AW-U).
\]
For $W<\Gamma_b$, the old lower-region row $W-\Gamma_b$ is negative,
but its multiplier is zero; thus no nonnegativity of this unused row
is needed. The upper row remains nonnegative because
$W\le1.25\Gamma_b$, and its increased multiplier is nonnegative by
\eqref{eq:qK}. Comparisons also remain nonnegative.
The same pointwise cancellation as in Lemma~\ref{lem:old} gives
$7.7547(A-A_b)1.25\Gamma_b=\Lambda_1$.

There is no gap at $W=\Gamma$, which belongs to the new low region.
The new first band ends at the unchanged upper endpoint; subsequent
bands and the tail are unchanged. The rule therefore covers every
possible selected-filter pattern exactly once, apart from the already
specified endpoint convention.
\end{proof}

\begin{remark}
No score function, filter-selection rule or continuous integral has
been replaced by a newly guessed estimate. The change is in which
existing certificate is used in the analysis and how its constants
are shifted. It does not infer positivity of an unused lower-region row.
\end{remark}

\section{Parameters, scalar certification and error budgets}\label{sec:admissible}
All terminating decimals below denote exact rationals. Here $\tau$ is the alias/checksum work parameter, $\eta n$ the large-prime bit scale, $\mu n$, $ln$ and $rn$ the sizes of the three alias blocks, $a$ the mask fraction within the large block, and $\gamma\mu n$ the small-prime bit scale. The constants $\eps$, $\zeta$ and $\theta$ specify pattern accuracy, preliminary-search slack and the bit-length guard. The frozen scalar certificate uses
\begin{center}\small
\begin{tabular}{@{}ll@{}}\toprule
Parameter & Value\\\midrule
$\tau$ & $0.48767228$\\
$\eta$ & $0.516133678389924$\\
$a$ & $0.008148495127574004$\\
$\mu$ & $0.4754356429062936$\\
$\gamma$ & $0.06005543072330252$\\
$l$ & $0.483775231690629$\\
$r=1-\mu-l$ & $0.0407891254030774$\\
$\eps,\zeta,\theta$ & $10^{-12}$ each\\
Common alias margin $s_*$ & $10^{-10}$\\
$m_{\rm fin}$ & $10^{-7}$\\
Advertised $\alpha$ & $0.4876723$\\\bottomrule
\end{tabular}
\end{center}
The exact budgets and internal total exponent are
\begin{align*}
 \Delta&=0.028461398389924,&\Gamma&=0.459210881610076,\\
 \beta=\tau+20\eps&=0.48767228002,&\alpha-\beta&=1.998\cdot10^{-8}.
\end{align*}
The construction parameters are changed, rather than only reducing the
rounding allowance in an unchanged theorem: $\tau<0.48767595$, the
internal parameter of~$\mathcal W$.

\begin{proposition}[Scalar feasibility]\label{prop:numeric}
The stated tuple satisfies every inequality in this section, with the
common alias margin $s_*=10^{-10}$, finite moment margin $10^{-7}$,
and error exponent $9.7468\cdot10^{-8}$. The continuous margin conclusion
is conditional on~\eqref{eq:oldmom}; the scalar comparisons themselves
use exact rational arithmetic with rigorous logarithm intervals.
\end{proposition}
\begin{proof}
The complete list of tested gaps and their definitions is in
Appendix~\ref{app:certificate}. In particular, all twenty alias/auxiliary
gaps exceed $s_*$, and $\zeta+\theta=2\cdot10^{-12}<s_*$.
The following decimal lower bounds have been rounded down from their
exact rational interval endpoints; their signs do not depend on a
floating-point tolerance.
\begin{center}\small
\begin{tabular}{@{}lr@{}}\toprule
Continuous margin family & Lower bound\\\midrule
Other 42 low certificates & $1.35967911\cdot10^{-7}$\\
Special filter-3 low certificate & $0.000949320934578$\\
All first bands, including extended filter 3 & $1.93876630\cdot10^{-6}$\\
Regions 2--8 & $3.8505054152$\\
Tail & $11.4443433093$\\\bottomrule
\end{tabular}
\end{center}
For clarity, the bounds are obtained from precisely these rational formulas:
\begin{align*}
 m_{\rm low}&=.00025+\eta-\eta_b-10.6428K,\\
 m_{3,0}&=.00025+\eta-\eta_b-\Lambda_0,\\
 m_{\rm first}&=.00025+\eta-\eta_b-8.4537(1.25)K,\\
 m_{2:8}&=4+\eta-\eta_b-418.36(15)K,\\
 m_9&=17+\eta-\eta_b-418.36(256)K/\Gamma_b.
\end{align*}
Every family bound exceeds $m_{\rm fin}=10^{-7}$.
Also $0<q<1$, $K>0$, $A<1/4$ and
$418.36(1+K/\Gamma_b)<419$.
Consequently the rounding-error exponent has the exact lower bound
\[
 m_{\rm fin}-(18+6\cdot419)\eps
 =10^{-7}-2532\cdot10^{-12}=9.7468\cdot10^{-8}>0.
\]
The logarithm enclosures and the exhaustive finite data checks supporting
these comparisons are proved/explained in Appendix~\ref{app:certificate}.
\end{proof}


\paragraph{Alias and size obligations.}
Write $r=1-\mu-l$ and $h(x)=-x\log_2x-(1-x)\log_2(1-x)$.
The eight conditions from Section~2.2 of~$\mathcal K$ are
\begin{align*}
 \gamma&<h(2a)/2, & l+(h(a)-\gamma)\mu&<\tau,\\
 r+(1-\gamma)\mu&<\tau,
 &1-\mu+(h(a)+h(1/2-a)-\gamma)\mu-\eta&<\tau,\\
 h(1/2-a)\mu+1/2-\tau&<\tau,
 &0<r<\tau,\quad l&<\tau,\\
 \mu(1-2a)&<l, &(\mu/2)\log_2 3&<\eta.
\end{align*}
The proof also uses $0<a<1/4$, $\mu(1-2a)<1/2$,
$h(2a)/2+h(a)<1-2a$, and $h(a)<h(1/2-a)$.
For $\delta=(\eta-(\mu/2)\log_2 3)/2$, require
\[
 \delta>1/2-\tau+\zeta,\qquad \delta>\theta,
\]
where $\theta$ is the input-guard exponent. These are the bad-half-block
and coverage conditions in A4/A5, stronger than $\delta>0$ alone.
The additional A7 work exponents are
\[
 \mu,\qquad r+1/2-\tau,\qquad1/2-\tau+\gamma\mu/2.
\]
All must lie below $\tau$ with a common positive work margin after
preliminary-search and guard costs. The exact checker must certify a
margin consistently below every relevant cost gap and place
$\zeta+\theta$ inside it, rather than checking only the displayed eight
inequalities. Prime generation also costs $2^{\eta n/2}\poly(n+b)$.

\paragraph{All 430 moments and approximation.}
For ordinary low certificates use loss $z_0K$; for the two special
certificates use~\eqref{eq:speciallow}--\eqref{eq:specialband}; for the
remaining finite bands and tail use Lemma~\ref{lem:old}. A per-certificate
new continuous margin is the inherited $m_i+(\eta-\eta_b)$ minus its loss.
One may bound ordinary low $z_0$ by $10.6428$ after removing filter~3,
first-band $z_0$ by $8.4537$, and other coefficients by $Z=418.36$,
as verified finite source data. The coarse tail bound is $B=256$
(Lemma~3 of~$\mathcal W$). Use these only with explicit positive
margins for every case, including the special low case.

For completeness, the score bound follows from eight coefficients of
absolute value at most2 and $1\le\kappa\le2$. Then $|r_v|,|s_v|\le16$.
Using the convolution notation $L=\log_2\mathcal P-\log_2(Z_fZ_g)$,
$B_*=\mathcal Y/\mathcal P-\log_2(Z_fZ_g)$ and
$H_*=\log_2Z_{f,\kappa}+\log_2Z_{g,\kappa}-\kappa\log_2(Z_fZ_g)$,
positivity and weighted averaging give $|L|,|B_*|\le64$ and
$|H_*|\le64\kappa$. Hence
$|w_v|=|(2-\kappa)B_*-2L+H_*|\le256$.
The only added multiplier is $z_0(A-A_b)$, so the sum is at most
$Z'=418.36(1+A-A_b)$.

Choose $m_{\rm fin}$ strictly smaller than the minimum new continuous
margin. Uniform finite-to-continuous convergence then supplies the
finite-circle bound $\eta-\log_2 I'_{\rm finite}>m_{\rm fin}$ for all
sufficiently large fixed prime cutoffs. Positivity and bounded denominators
of the fixed filters give an effective cutoff as in Lemma~3.2
of~\cite{openai}; this is an inherited analytical argument, not a new
numerical integration.
The prescribed score errors are $3\eps$ for $U$ and $4\eps$ for $W$.
If $A<1/4$, every modified affine row differs from its coordinate average
by at most $6\eps$, including comparisons; shifted constants add no error.
Thus ideal domination gives
\[
 -\alpha U\le-\Delta+\frac1n\sum_{j=1}^n z'\cdot G'(x_j)+6Z'\eps.
\]
The inherited sampled-pattern mean-square bound has decay exponent
$m_{\rm fin}-(18+6Z')\eps$, which must be strictly positive.
The $2^{O(\sqrt n)}$ patterns and polynomially many transformed targets
are absorbed by this decay. The original good-table argument and finite
precision then give the same constant absolute-error guarantees.

\section{Correctness, worst-history cost and exponent absorption}
\begin{proposition}[Conditional algorithmic implication]\label{prop:algorithm}
Assume the inherited hypotheses stated in Section~1, and suppose the
fixed rational parameters satisfy all conditions and strict margins
in Section~\ref{sec:admissible}. The modified construction has
\[
 T(n,b)\le2^{(\tau+20\eps)n+O(n^{0.7})}\poly(n+b+2)
\]
on every guarded random execution. For every fixed $c$ and any fixed
$\alpha>\tau+20\eps$, its full program has time
$O_c(2^{\alpha n})$ on $b\le n^c$, halts on every input, and has
per-input correctness at least $2/3$ after a fixed small-input fallback.
\end{proposition}
\begin{proof}
Here LD is the inherited low-deficiency condition for the transformed
weights $c_i=Ca_i+u_i$, with the tie weights $u_i$ fixed:
\[
 \begin{gathered}
 \Sigma_c(V)=\Bigl\{\sum_{i\in I}c_i:I\subseteq V\Bigr\},\qquad
 d_c(V)=|V|-\log_2|\Sigma_c(V)|,\\
 \mathbb E_{V\sim\operatorname{Unif}\binom{[n]}{\lfloor n/2\rfloor}}
 d_c(V)\le(\tfrac12-\tau+\zeta)n.
 \end{gathered}\tag{LD}
\]
The expectation is over all half-blocks of the indicated size, not over
the tie weights or the observed preliminary-search outcome. The inherited
transformation takes $N_0$ to be the power of two in $[12n,24n)$,
$u_i\in\{1,\ldots,N_0\}$ independently uniform before fixing them,
$C=1+nN_0$, and transformed targets $Ct+s$ for $0\le s\le nN_0$,
where $a_i$ and $t$ are the original input weights and target.
This is precisely Kapila's condition (LD) in Section~4, with his $\tau_0$
written as our $\tau$; see also Section~2.1 and Appendix~B.1
of~\cite{kapila}. It is the existing case distinction in the inherited
argument, not an additional assumption on all inputs.
The preliminary search either gives an exact answer or continues.
For each fixed choice of tie weights the analysis divides according
to LD, not according to the observed continuation event. Draw the main
procedure's independent randomness even if an earlier exact answer
makes it unused. The inherited preliminary guarantee covers failure
of LD. Under LD, the source alias proof uses the checked coverage,
size and work inequalities; caps limit every run, and its expected-work
estimate makes cap failure $o(1)$. The modified checksum domination
and moments were proved above; the inherited table/estimator analysis
therefore has error $o(1)$ uniformly in the guarded fixed input.

Use the same retained phases in the modular checksum and alias
subtraction. Complete, deduplicated alias phase sums cancel all
modular representations whose unreduced sum is different from the
target. The remaining exact bin is zero on a NO input. On a YES input,
isolation supplies a unique exact bin with probability at least $11/12$,
and that bin has modulus1. A combined checksum/alias arithmetic error
less than $1/4$ makes threshold $1/2$ correct on these events.
All additional failure probabilities are $o(1)$ uniformly above a fixed
cutoff; choosing their sum below $1/12$ gives success at least $5/6$,
hence $2/3$.

All samplers, emitted lists, visits and expansions have deterministic
caps, checked before the next write. Random-prime generation by bounded
trial division costs $2^{\eta n/2}\poly(n+b)$ and is below the certified
work budget. Retained randomness and polynomial-bit rational arithmetic
are charged, as in the inherited implementation. The cap exponents,
score accuracy and block-pattern factors give the displayed bound;
none is replaced by an expected running time. Required constants and
fixed tables are generated by the program, not supplied as advice.

Use the exact meet-in-the-middle fallback below a fixed cutoff and
when the integer guard fails. For integer $\theta^{-1}$ the guard is
$\theta^{-1}\lceil\log_2(b+1)\rceil\le n$; it implies
$b+2\le2^{\theta n+1}$, as required in alias collision estimates.
For each fixed $c$, every $b\le n^c$ eventually satisfies this guard.
All cutoffs and algorithm parameters are independent of $c$; only the
final asymptotic comparison may depend on it. If $\beta=\tau+20\eps$,
then $O(n^{0.7})+\log_2\poly(n+n^c+2)=o_c(n)$.
For fixed $\alpha>\beta$ this is eventually at most $(\alpha-\beta)n$,
proving ordinary $O_c(2^{\alpha n})$. The inherited address-space proof
stores multiword operands and counts each allocated word; its logarithmic
storage bound fits $\lceil4(n+b+\log_2(n+2))\rceil$-bit addresses.
\end{proof}

\section{Verification scope, attribution and limits}
Floating-point discovery is not certification. The certificate fixes rational parameters, encloses entropy values with justified error, checks every numerical inequality above and the complete finite source transcription, and supplies one small verification command that does not rerun the optimizer (Appendix~\ref{app:certificate}). Continuous moment inequalities remain
the explicit inherited hypotheses~\eqref{eq:oldmom}; a scalar certificate
does not establish them or the whole upstream algorithm. No claim of
framework optimality follows from a failed search.

The current campaign uses no locally verified Lean build: no Lean/elan
installation was available in the inspected environment. Wharton's
later formalization has a different, documented finite-data scope, and
its continuous moment hypotheses do not disappear. The live verifier
has no trusted Subset Sum statement at startup~\cite{rules}.
Local certificate validity, manuscript proof completeness, and site
verification are distinct claims.

\paragraph{Acknowledgments and provenance.}
The intended author is Yanni Bills. The research team used AI assistance
for source investigation/exposition (Researcher-Ed\_Lorenz), bounded
parameter discovery and certificate implementation (Computational-Richardson),
coordination/integration (Coordinator), and independent criticism
(Reviewer-Procopius). The inherited construction and proofs are credited
to OpenAI and Annanay Kapila, and the original scalar transfer and exact
numerical machinery to Joshua Wharton. The new boundary treatment must
be assessed separately from those inherited results. Source commits,
unchanged downloads and verification scope are recorded with the repository.
\appendix
\section{Exact certificate, gap definitions and verification scope}\label{app:certificate}
\paragraph{Logarithm enclosures.}
For $1\le x\le2$, put $z=(x-1)/(x+1)\in[0,1/3]$. Integrating the
geometric series for $1/(1-z^2)$ gives
\[
 \ln x=2\sum_{j=0}^{N-1}\frac{z^{2j+1}}{2j+1}+R_N,
 \qquad
 0\le R_N\le\frac{2z^{2N+1}}{(2N+1)(1-z^2)}.
\]
The tail bound follows by replacing every remaining denominator by
$2N+1$ and summing the resulting geometric series. For positive rational
$x$, multiply or divide by powers of2 to write $x=2^k y$ with $1\le y<2$.
Then $\log_2x=k+\ln y/\ln2$. Use the same series with $N=48$ to enclose
both logarithms; the lower endpoint for $\ln2$ is positive.
Interval division therefore gives a rational enclosure for every
$\log_2x$. The entropy bounds follow by interval arithmetic in
$-x\log_2x-(1-x)\log_2(1-x)$. No floating-point number is used to decide
an inequality. Addition and multiplication use the extremal endpoint
combinations; division is used only on intervals not containing zero.
These operations preserve enclosure by elementary order properties.

\paragraph{Exhaustive scalar gap definitions.}
Let $H=h(1/2-a)$ and $d=(\eta-(\mu/2)\log_2 3)/2$.
Each row below is a positive gap; the reported lower bounds are conservative
roundings, not the values used for exact decisions. These are all the
alias/auxiliary gap families checked, including those hidden in the source's
``far below'' estimates.
\begin{center}\scriptsize
\begin{tabular}{@{}p{.70\linewidth}r@{}}\toprule
Gap & Lower bound\\\midrule
$h(2a)/2-\gamma$ & $9.99999\cdot10^{-10}$\\
$\tau-l-(h(a)-\gamma)\mu$ & $4.00657\cdot10^{-9}$\\
$\tau-r-(1-\gamma)\mu$ & $4.00657\cdot10^{-9}$\\
$\tau-1+\mu-(h(a)+H-\gamma)\mu+\eta$ & $4.00657\cdot10^{-9}$\\
$2\tau-1/2-H\mu$ & $7.01315\cdot10^{-9}$\\
$\tau-l$ & $0.00389704$\\
$\tau-r$ & $0.44688315$\\
$r$ & $0.04078912$\\
$l-\mu(1-2a)$ & $0.01608775$\\
$\eta-(\mu/2)\log_2 3$ & $0.13935984$\\
$\tau-\mu$ & $0.01223663$\\
$2\tau-1/2-r-\zeta$ & $0.43455543$\\
$2\tau-1/2-\zeta-\gamma\mu/2$ & $0.46106831$\\
$\tau-\eta/2$ & $0.22960544$\\
$1/2-\mu(1-2a)$ & $0.03231252$\\
$1-2a-h(2a)/2-h(a)$ & $0.85539536$\\
$H-h(a)$ & $0.93155619$\\
$d-\theta$ & $0.06967992$\\
$d-(1/2-\tau+\zeta)$ & $0.05735220$\\
$0.013-(1/2-\tau+\zeta)$ & $0.00067227$\\\bottomrule
\end{tabular}
\end{center}
All exceed $s_*$. The work inequalities without an explicit $\zeta$ in
the table have at least $s_*$ before spending preliminary-search slack;
the sorting/prime inequalities listed with $\zeta$ already have that
slack charged. Charging it once more everywhere is harmless because
$\zeta+\theta<s_*$ and leaves a positive margin. Rounding of subset
sizes contributes $O(1)$ or $O(\log n)$ terms, dominated by those strict
linear gaps. The small-layer entropy inequality implies the source's
binomial size condition with exponential slack, and its complementary
sizes use $0<a<1/4$ and $\mu(1-2a)<\min(l,1/2)$.
Thus the concentration/coverage estimates in A3--A6 and every A7 work
item retain their prescribed asymptotic inequalities, not only C1--C8.

\paragraph{Finite source data and complete region coverage.}
The original baseline TeX has SHA256
\begin{center}\scriptsize
\texttt{b397531cc5997e88c439d5290774861b46a34aa32d9cd71fa7dcced130359298}.
\end{center}
The checker requires exactly filter indices0--42 and all ten regions for
each, rejecting missing or duplicate entries. It checks sixteen rational
coefficients per filter, each of absolute value at most2, and
$1\le\kappa\le2$. Each certificate has the required number of multipliers,
nonnegative entries with sum at most418.36, valid distinct comparison
indices excluding its own filter, and the source's domination condition
$z_0\ge1$ or $z_0L\ge1$. It authenticates the special rows used in
Lemma~\ref{lem:new}, verifies the zero lower multiplier exactly, and
checks the maxima10.6428 and8.4537 used in the family estimates.
The complete $43\cdot10$ index product, the explicit partition in
Lemma~\ref{lem:new} and the unchanged remaining boundaries establish
coverage of all selected-filter patterns; a floating sample of regions
would not suffice. Numerical source transcription does not prove any
continuous moment integral in~\eqref{eq:oldmom}.

\paragraph{Compact verification and discovery separation.}
The independently written Python checker is
\texttt{numerics/verify\_certificate.py}; exact inputs are
\texttt{numerics/data/candidate.json}. Given the unchanged pinned
baseline, run from the delivery repository
\begin{quote}\small\ttfamily
python numerics/verify\_certificate.py --baseline /path/to/PROOF\_0.4877.tex
\end{quote}
The baseline remains an explicit external input rather than an unlicensed
full manuscript redistributed in the repository. Its URL, revision and
hash are recorded. The checker needs Python's standard library and does
not run the optimizer or require a numerical integration library.
The separate discovery batch used 24 floating optimizer starts with a
recorded batch time of approximately $0.447$ seconds; this is parameter-search
time, not a Subset Sum solver runtime, and its success is not evidence for
the inequalities.
The exact checker result and independent review are separate recorded
outputs. Negative tests must reject broken parameters, changed special
rows and incomplete source coverage; their precise executed scope is
reported with the code. No end-to-end Subset Sum implementation was run.

\begingroup\small
\bibliographystyle{plain}
\bibliography{references}
\endgroup
\end{document}
